\section{Conclusion}
\label{sec:conclusion}

% src: Sections 5-8 (numbers as cited there); CLAIM_EVIDENCE_MAP.md C1-C9
In conclusion, this work has shown, in three simulated single-BS studies, that using the history of a user's reports, rather than the current report alone, consistently lowers the error of the best model across propagation settings, BS placements and model families. In the macro cell, five steps of history lower the error of devices that report an AoA by $15$ to \SI{21}{\percent} in all seven model families tested, so that the input matters far more than the estimator, and LOS propagation, interfering BSs and a blocked line of sight leave that result unchanged. How the history is used still matters: on the controlled grid and in the $25\times25$ runs, sequence models use windows of fifteen reports where tree ensembles stop at about three, the order of the reports carries part of that gain, and a motion-model filter, which suits the regular grid walk, uses history better still. What history reads differs between the scales; in the macro cell the gain is the averaging of noisy angle reports, chiefly elevation, and it shrinks to $2$ to \SI{7}{\percent} when the angle error persists across reports. Finally, the gain depends on the position information a single report carries and on whether it is shared between users: single-antenna devices in the uniform-NLOS macro cell, whose reports carry nothing that transfers to new users, gain no more than about \SI{1}{\percent}, and the grid's gain without an angle rests on a channel map learned for each user, whereas in an environment shared by all users, history lowers the error of unseen single-antenna devices by about $6$ to \SI{7}{\meter} with a recurrent model, under uniform NLOS as well.

Four follow-ups would settle what remains. Measured AoA traces would show how correlated real angle errors are, which sets where the gain lies between the independent and persistent cases. Simulations with variable step length or sampling rate would separate the number of reports from the distance they cover. A direct examination of what the tree ensembles infer from RSS history under mixed propagation would test the account offered for the single-antenna gain. Finally, a fully shared environment, including the small-scale fading, and the same test on the controlled grid would show whether the within-user gain carries over to new users there as it does, in part, in the partly shared macro-cell environments.
